\documentclass[12pt]{article}
\newcommand{\DocShort}{Tasa decreciente}
\input{preamble_population.tex}
\begin{document}
\doccover{Método de la tasa decreciente de crecimiento}{2026-05}
\handoutintro{Proyectar la población reduciendo en cada período censal la tasa porcentual de crecimiento observada más recientemente.}{Ciudades maduras cuyo ritmo de crecimiento empieza a desacelerarse.}

\section{Datos ejemplo}
Se usará una serie con períodos censales de $10$ años; por tanto, las tasas porcentuales del ejemplo son tasas por década.

\begin{center}
\begin{tabular}{lcccccc}
\toprule
Año & 1970 & 1980 & 1990 & 2000 & 2010 & 2020 \\
\midrule
Población & 145,800 & 170,600 & 197,100 & 223,200 & 247,000 & 267,800 \\
\bottomrule
\end{tabular}
\end{center}

\section{Ecuaciones mínimas}
La tasa porcentual de crecimiento del período censal $i$ es:
\begin{equation}
\%\Delta_i=\frac{P_i-P_{i-1}}{P_{i-1}}\times 100
\end{equation}
donde $P_{i-1}$ es la población al inicio del período y $P_i$ es la población al final del período.

La disminución media de la tasa porcentual se obtiene como:
\begin{equation}
d=\frac{1}{M}\sum_{j=1}^{M} d_j
\end{equation}
donde $d_j$ es cada disminución reciente de la tasa porcentual y $M$ es el número de disminuciones consideradas.

La primera tasa futura corregida es:
\begin{equation}
\%\Delta_{\text{futura,1}}=\%\Delta_{\text{última}}-d
\end{equation}

\weightednote{Si los intervalos censales no son regulares, conviene estandarizar las tasas a una base anual antes de promediarlas y luego usar un promedio ponderado con pesos iguales a la duración de cada período en años.}

\section{Procedimiento}
\begin{enumerate}
\item Calcular los incrementos absolutos.
\item Calcular las tasas porcentuales de crecimiento por período censal.
\item Calcular las disminuciones entre tasas consecutivas.
\item Promediar las disminuciones recientes.
\item Corregir la última tasa observada y proyectar período por período.
\end{enumerate}

\section{Ejemplo resuelto}
\begin{center}
\begin{tabular}{crrrr}
\toprule
Año & Población & \makecell{Incremento\\del período} & \makecell{\% de crecimiento\\por década} & \makecell{Disminución\\de tasa} \\
\midrule
1970 & 145,800 & -- & -- & -- \\
1980 & 170,600 & 24,800 & 17.01 & -- \\
1990 & 197,100 & 26,500 & 15.53 & 1.48 \\
2000 & 223,200 & 26,100 & 13.24 & 2.29 \\
2010 & 247,000 & 23,800 & 10.66 & 2.58 \\
2020 & 267,800 & 20,800 & 8.42 & 2.24 \\
\bottomrule
\end{tabular}
\end{center}

Tomando las tres disminuciones más recientes:
\[
d=\frac{2.29+2.58+2.24}{3}=2.37\ \text{puntos porcentuales por década}
\]

La última tasa observada fue $8.42\%$ por década. Entonces:
\[
\%\Delta_{2030}=8.42-2.37=6.05\%
\]
\[
P_{2030}=\num{267800}(1+0.0605)=\num{284003}
\]
\[
\%\Delta_{2040}=6.05-2.37=3.68\%, \qquad P_{2040}=\num{294453}
\]
\[
\%\Delta_{2050}=3.68-2.37=1.31\%, \qquad P_{2050}=\num{298306}
\]

\resultbox{La población proyectada es 284,003 habitantes en 2030, 294,453 en 2040 y 298,306 en 2050.}

\section{Teoría breve}
Este método intenta representar la parte final de una curva en S: la ciudad sigue creciendo, pero lo hace cada vez a menor ritmo.

\section{Observaciones de uso}
\begin{itemize}
\item No conviene si la ciudad todavía muestra aceleración.
\item Debe distinguirse entre tasa porcentual por período censal y disminución media de la tasa por período censal.
\item La tasa futura no debe hacerse negativa sin justificar un escenario de decrecimiento real.
\end{itemize}
\end{document}
